\documentclass[11pt]{article}

\usepackage[margin=1in]{geometry}
\usepackage{amsmath,amssymb,amsthm}
\usepackage{mathtools}
\usepackage{enumitem}
\usepackage{booktabs}
\usepackage{hyperref}
\usepackage{xurl}
\usepackage{microtype}
% BibLaTeX in-file (single-file paper)

\newcommand{\takeaway}[1]{\paragraph{Takeaway.}\emph{#1}}

\title{Calibration Recipes for Anonymous DHT Deployments:\\ Budgets, Horizons, and Alarm Tax}
\author{}
\date{February 24, 2026}

\begin{document}
\maketitle

\begin{abstract}
To keep the main Anonymous DHT papers referee-tight and non-redundant, we collect the ``numbers-first'' calibration tables and rule-of-thumb conversions that otherwise get repeated:
(i) witness-amplification horizons for state-dependent leakage,
(ii) translating observation cadence into time-to-detection,
(iii) the inevitable alert-oracle tax for closed-view auditing.
We also include pointers to the CPPC manifest interface (AnonDHT~State~3) without reprinting schemas or cards.\cite{state3}
Core definitions and theorems live in the State series and Synthesis notes; this paper is purely an operational bridge.
\end{abstract}
\paragraph{Scope.}
This is an operational bridge: ``numbers-first'' recipes and tables that help size witness cadences, horizons, and alarm costs without repeating algebra across the core papers.

\paragraph{Units.} Budgets are published in bits (use $\log_2$ and $2^{\epsilon}$); results stated in nats can be converted by dividing by $\ln 2$ (see Synthesis~8).\cite{series:notation}

\paragraph{Imports.}
We import the threat/window checklist and trace-interface vocabulary.\cite{series:threatwindows,series:traceifaces}
Amplification bounds are imported from the State series (Papers~1 and~3), and identification sizing rules from the endpoint bridge note.\cite{state1,state3,synth5}

\paragraph{Exports.}
We export compact calibrations: (i) quadratic horizon scaling for persistent witness gaps, (ii) cadence-to-horizon translations, and (iii) a tabulated alert-oracle ``tax'' lower bound for closed-view auditing.
For the CPPC manifest interface, see AnonDHT~State~3.\cite{state3}

\paragraph{Artifact policy.}
This archive is source-only; the full CPPC schema and supporting tooling are available on request.\cite{series:artifactpolicy}


\section{Series context and what this note owns}
We use the threat/window checklist \cite{series:threatwindows} and the trace-interface vocabulary \cite{series:traceifaces}.
This note owns only ``back-of-the-envelope'' translations that help readers size cadence, windows, and budgets without re-deriving the same algebra in every paper.

\section{Recipe 1: witness amplification needs quadratic horizons}
State series Paper~1 \cite{state1} shows that even a tiny, persistent witness bias amplifies under repetition.
A convenient calibration is the Hoeffding-style rule of thumb
\[
T \approx \frac{2\log_2(1/\eta)}{\Delta^2},
\]
where $\Delta$ is the per-epoch witness gap and $\eta$ is a target total error probability.

\begin{table}[h]
\centering
\begin{tabular}{@{}rrrr@{}}
\toprule
Per-epoch gap $\Delta$ & Target total error $\eta$ & Rule-of-thumb $T$ & Interpretation \\
\midrule
0.01 & 0.05 & $\approx 6\times 10^4$ & tiny bias needs weeks at minute cadence (or days at 10s cadence) \\
0.02 & 0.05 & $\approx 1.5\times 10^4$ & $\sim$10 days at minute cadence (or $\sim$2 months at 5min cadence) \\
0.05 & 0.05 & $\approx 2.4\times 10^3$ & $\sim$40 minutes at 1s cadence (or $\sim$7 hours at 10s cadence) \\
0.10 & 0.05 & $\approx 6\times 10^2$ & $\sim$10 minutes at 1s cadence (or $\sim$1.2 hours at 10s cadence) \\
\bottomrule
\end{tabular}
\caption{Calibration for witness amplification (Paper~1 \cite{state1}). Epoch vocabularies vary across systems (e.g.\ Nym uses minute-scale epochs \cite{nym_cryptoecon_reward_sharing}); the key scaling is the quadratic dependence on $\Delta$.}
\label{tab:calibration}
\end{table}

\takeaway{If an adversary can extract even a 1--2\% per-epoch witness gap and the system runs at high cadence, ``state'' becomes a fast privacy amplifier unless it is refreshed/privatized.}

\section{Recipe 2: budget-to-horizon translation (cadence \texorpdfstring{$\rightarrow$}{->} time)}
If a monitoring/evidence stream yields $T$ independent-ish epochs per window, the same $T$ can correspond to very different wall-clock horizons depending on cadence and sampling.
Table~\ref{tab:budget-horizon} gives an illustrative translation assuming the rule of thumb above with $\eta=0.05$.

\begin{table}[h]
\centering
\begin{tabular}{@{}rrrr@{}}
\toprule
Per-epoch gap $\Delta$ & $T$ (epochs) & Horizon at $10^3$/day & Horizon at $10^5$/day \\
\midrule
0.01 & $6\times 10^4$ & $\approx 60$ days & $\approx 14$ hours \\
0.02 & $1.5\times 10^4$ & $\approx 15$ days & $\approx 3.6$ hours \\
0.05 & $2.4\times 10^3$ & $\approx 2.4$ days & $\approx 35$ minutes \\
\bottomrule
\end{tabular}
\caption{Budget-to-horizon calibration for witness amplification (Paper~1 \cite{state1}). If an observer sees only a fraction $q$ of epochs, multiply horizons by $1/q$. If state is refreshed every $W$ days, amplification effectively saturates once $T$ exceeds the refresh horizon.}
\label{tab:budget-horizon}
\end{table}

\section{Recipe 3: the alert-oracle tax is logarithmic but not free}
Closed-view auditing (State series Paper~3 \cite{state3}) can require ``alarms'' or public alerts. If an alert has miss probability $\beta$ and false-alarm probability $\alpha$, the paper's alert-oracle bound yields a necessary privacy cost roughly
\[
\varepsilon \gtrsim \log_2\!\left(\frac{1-\beta}{\alpha}\right),
\]
independent of mechanism details.

\begin{table}[h]
\centering
\begin{tabular}{@{}rrr@{}}
\toprule
Miss prob.\ $\beta$ & False-alarm $\alpha$ & Lower bound $\varepsilon \gtrsim \log_2((1-\beta)/\alpha)$ \\
\midrule
0.10 & 0.05  & 2.89 \\
0.10 & 0.01  & 4.50 \\
0.10 & 0.001 & 6.80 \\
0.05 & 0.05  & 2.94 \\
0.05 & 0.01  & 4.55 \\
0.05 & 0.001 & 6.86 \\
0.01 & 0.05  & 2.99 \\
0.01 & 0.01  & 4.60 \\
0.01 & 0.001 & 6.90 \\
\bottomrule
\end{tabular}
\caption{Calibration for the alert-oracle tax (Paper~3 \cite{state3}). Tight alarms cost nontrivial $\varepsilon$; pushing $\alpha$ down by 1000x costs about 6.9 nats $\approx 10$ bits.}
\label{tab:alert-tax-calibration}
\end{table}

\takeaway{Alarm quality is a privacy knob: to get rare false alarms you must usually spend a meaningful privacy budget, even before any ``payload'' disclosures.}

\section{Recipe 4: endpoint sizing via MaxL and odds inflation (pointer)}
For identification and shortlisting conversions (e.g.\ ``what MaxL budget prevents exact identification among $K$ candidates?''), we do not restate algebra here.
Use the endpoint bridge note \cite{synth5}, which relates realized odds inflation, pointwise maximal leakage, and maximal leakage and records the standard $\log_2(\alpha K)$ sizing rule.

For destination-privacy equalization of routing transcripts, MC-EQ/CoverSketch provides an explicit session scaling bound under uniform covers: $\mathcal{L}_{\max}\le T\log_2(1+2k\delta_{eq})$ (receipt-friendly and conservative).\cite{series:mceq}

\section{CPPC schema pointer}
For the CPPC manifest interface (schema excerpt, example cards, and signing/logging guidance), see the closed-view auditing note (AnonDHT~State~3).\cite{state3} The full schema and tooling are distributed out-of-band on request.\cite{series:artifactpolicy}

\section{Limitations and open problems}
These recipes ignore correlation structure, adaptivity, and adversarial selection effects; they are intended for sanity checks, not tight privacy accounting.
Where correlations matter (e.g.\ retries in the Boss Fight), prefer mechanism-specific bounds and then convert to endpoint budgets via \cite{synth5}.

\section{Takeaways}
\begin{itemize}[leftmargin=*]
\item \textbf{Quadratic amplification:} persistent per-epoch witness gaps ($\Delta$) translate to horizons ($T$) on the order of $\log_2(1/\eta)/\Delta^2$ (Recipe~1).
\item \textbf{Cadence is a privacy knob:} the same $T$ epochs can mean hours or months depending on sampling rate; partial visibility scales horizons by $1/q$ (Recipe~2).
\item \textbf{Alarms cost budget:} tightening false-alarm probability pushes a nontrivial alert-oracle lower bound (Recipe~3); plan to spend privacy budget on operational signals, not just payloads.
\item \textbf{Don't re-derive endpoint sizing:} convert mechanism bounds to endpoint identification/shortlisting budgets via the endpoint bridge note (Recipe~4).
\end{itemize}

\begin{thebibliography}{99}\small

\bibitem{state1}
Anonymous.
\newblock State-Dependent Anonymity in Anonymous DHTs: Witnesses, Amplification, and Stability.
\newblock 2026.
\newblock AnonDHT State series Paper~1 (this archive).

\bibitem{state3}
Anonymous.
\newblock Closed-View Auditing for Stateful Anonymity: Odometers, Alarms, and Machine-Checkable CPPCs.
\newblock 2026.
\newblock AnonDHT State series Paper~3 (this archive).

\bibitem{synth5}
Anonymous.
\newblock Endpoint Metrics Bridge for Anonymous Overlays: Odds Inflation, PML, and MaxL.
\newblock 2026.
\newblock Synthesis series note (Synthesis~5) in this archive.

\bibitem{series:traceifaces}
Anonymous.
\newblock Trace Interfaces for Anonymous DHTs: Observation Tiers, Committee-Contact Traces, and Exposure Normal Forms.
\newblock 2026.
\newblock Synthesis series note (Synthesis~3) in this archive.

\bibitem{series:threatwindows}
Anonymous.
\newblock Threat Models and Repetition Windows for Anonymous DHTs.
\newblock 2026.
\newblock Synthesis series note (Synthesis~4) in this archive.

\bibitem{nym_cryptoecon_reward_sharing}
Xinmu Alexis Cao and others.
\newblock Nym Cryptoeconomic Paper (reward sharing and epochs).
\newblock 2021.
\newblock Public Nym document; used only as an example of an epoch/cadence vocabulary.


\bibitem{series:mceq}
Anonymous.
\newblock MC-EQ and CoverSketch: Destination Privacy via Markov-Chain Equalization in Anonymous DHT Lookups.
\newblock 2026.
\newblock Release \& destination Paper~B (this archive).

\bibitem{series:artifactpolicy}
Anonymous.
\newblock Artifact policy (source-only archive; artifacts available on request).
\newblock Series note, 2026.

\bibitem{series:notation}
Anonymous.
\newblock Notation and Minimal Model for the One-True-Archive.
\newblock Series synthesis note (Synthesis~8), 2026. (This archive.)

\end{thebibliography}

\end{document}
